\documentclass[11pt,a4paper]{article}
% Shared preamble for RuPhO English problem statements (match T1 2026 style).
% Use: \input{latex/rupho_en_preamble.tex} after \documentclass.
\usepackage[margin=1in]{geometry}
\usepackage{amsmath,amssymb}
\usepackage{graphicx}
\usepackage{float}
\usepackage{enumitem}
\usepackage{microtype}
\usepackage{caption}
\usepackage{tabularx}
\usepackage{hyperref}

\captionsetup{font=small,labelfont=bf,skip=6pt}

\setlist[itemize]{leftmargin=1.2cm}
\setlist[enumerate]{leftmargin=1.2cm}

% One outer box per subproblem: [id | statement | points] (no inner rules)
\newcommand{\problembox}[3]{%
  \par\addvspace{0.42em}%
  \noindent
  \setlength{\fboxsep}{7.5pt}%
  \setlength{\fboxrule}{0.95pt}%
  \fbox{%
    \begin{minipage}{\dimexpr\linewidth-2\fboxsep-2\fboxrule\relax}
    \begin{tabularx}{\linewidth}{@{}>{\raggedright\arraybackslash\bfseries}p{2.1em} @{\hspace{0.9em}} >{\raggedright\arraybackslash}X @{\hspace{0.9em}} >{\raggedleft\arraybackslash\bfseries}p{2.4em}@{}}
    #1 & #2 & #3
    \end{tabularx}
    \end{minipage}%
  }%
  \par\addvspace{0.32em}%
}


\title{\textbf{Down the straight line that leads you up}}
\author{\normalsize X16 (2016) --- English translation}
\date{}
\begin{document}
\maketitle

An ideal polyatomic gas participates in a process, the states of which are depicted in the figure by points on the straight line $p=p_M (1- \alpha V)$, where $\alpha$ is a certain constant coefficient, $V$ is the current value of the volume. Molar heat capacity $C_V$ does not depend on temperature.

\begin{figure}[H]
    \centering
    \includegraphics[width=0.72\textwidth]{T5_EN_assets/fig1.jpg}
    \caption{Figure 1.}
\end{figure}

\problembox{A1}{At what pressure $p_0$ does the heat capacity of the gas become zero?}{2.00}

\problembox{A2}{Now suppose the process involves an ideal Clausius gas, whose equation of state is
\[
p=\frac{RT}{V-b}.
\]
For which gas (ideal polyatomic or ideal Clausius) is the temperature at the point of maximum heat capacity higher? Calculate the temperature difference $\Delta T$ between those points. Given:
$b = 44~\text{cm}^3/\text{mol}$, $R = 8.314~\text{J}/(\text{mol}\cdot\text{K})$, $p_M = 1.51~\text{MPa}$, and $b \ll V_M$, where $V_M=p_M/\alpha$ is a fixed volume.}{3.00}

\problembox{A3}{What value does the molar heat capacity $C$ of an ideal polyatomic gas participating in the process described in paragraph A1 tend to as the volume tends to zero?}{2.00}

\vspace{1em}
\noindent\textit{Source:} \url{https://pho.rs/p/4003}

\end{document}